\documentclass{article}
\usepackage[T1]{fontenc}
\usepackage{lmodern}
\usepackage{graphicx}
\usepackage{amsmath,amssymb}
\usepackage[a4paper,margin=2cm]{geometry}
\usepackage[absolute,overlay]{textpos}
\setlength{\TPHorizModule}{1mm}
\setlength{\TPVertModule}{1mm}
\usepackage[indonesian]{babel}
\usepackage{hyperref}
\setlength{\emergencystretch}{2em}
% unit: Halaman 26

\begin{document}

% === Halaman 26 (llm) ===

\section*{Cipher Block Chaining(CBC)}

\begin{itemize}
    \item Tujuan: membuat ketergantungan antar blok (mengatasi kelemahan mode ECB).
    \item Setiap blok cipherteks bergantung tidak hanya pada blok plainteksnya tetapi juga pada seluruh blok plainteks sebelumnya.
    \item Hasil enkripsi blok sebelumnya di-umpan-balikkan ke dalam enkripsi blok yang current.
    \item Enkripsi blok pertama memerlukan blok semu ($C_0$) yang disebut IV (initialization vector).
    \item IV dapat diberikan oleh pengguna atau dibangkitkan secara acak oleh program.
    \item Pada dekripsi, blok plainteks diperoleh dengan cara meng-XOR-kan IV dengan hasil dekripsi terhadap blok cipherteks pertama.
\end{itemize}

\end{document}
